\section{模型求解与数据测量}

\subsection{不同城市的健壮性测量}

我们用蒙特卡洛算法，在每个破坏节点百分比随机抽样100组破坏节点组合，然后求取平均瞬时健壮性，再对此进行积分计算，并画出城市健壮性曲线。如图10：

我们定义健壮性性能（即被积函数）曲线斜率变化最大处即健壮性显著变化处。我们的测量结果如表2所示
\begin{figure}[t]
  \centering
  \begin{subfigure}[t]{0.24\linewidth}
    \includegraphics[width=\linewidth]{figures/Q2/Chengdu_Robustness_Random_BC.pdf}
    \caption{Chengdu Robustness}
    \label{fig:cache}
  \end{subfigure}\hfill
  \begin{subfigure}[t]{0.24\linewidth}
    \includegraphics[width=\linewidth]{figures/Q2/Dalian_Robustness_Random_BC.pdf}
    \caption{Dalian Robustness}
    \label{fig:scores}
  \end{subfigure}
  \begin{subfigure}[t]{0.24\linewidth}
    \includegraphics[width=\linewidth]{figures/Q2/Dongguan_Robustness_Random_BC.pdf}
    \caption{Dongguan Robustness}
    \label{fig:cache}
  \end{subfigure}\hfill
  \begin{subfigure}[t]{0.24\linewidth}
    \includegraphics[width=\linewidth]{figures/Q2/Harbin_Robustness_Random_BC.pdf}
    \caption{Harbin Robustness}
    \label{fig:scores}
  \end{subfigure}
  \begin{subfigure}[t]{0.24\linewidth}
    \includegraphics[width=\linewidth]{figures/Q2/Qingdao_Robustness_Random_BC.pdf}
    \caption{Qingdao Robustness}
    \label{fig:cache}
  \end{subfigure}\hfill
  \begin{subfigure}[t]{0.24\linewidth}
    \includegraphics[width=\linewidth]{figures/Q2/Quanzhou_Robustness_Random_BC.pdf}
    \caption{Quanzhou Robustness}
    \label{fig:scores}
  \end{subfigure}
  \begin{subfigure}[t]{0.24\linewidth}
    \includegraphics[width=\linewidth]{figures/Q2/Shenyang_Robustness_Random_BC.pdf}
    \caption{Shenyang Robustness}
    \label{fig:cache}
  \end{subfigure}\hfill
  \begin{subfigure}[t]{0.24\linewidth}
    \includegraphics[width=\linewidth]{figures/Q2/Zhengzhou_Robustness_Random_BC.pdf}
    \caption{Zhengzhou Robustness}
    \label{fig:scores}
  \end{subfigure}
  \caption{不同城市的健壮性积分曲线}
\end{figure}

\begin{table}[htbp]
  \centering
  \caption{各城市公交网络鲁棒性综合评估（Q2）}
  \label{tab:Q2_robustness_bc}
  \begin{tabular}{lcc}
    \toprule
    \textbf{城市} & \(R_{bc}\)（健壮性积分） & \(q\)（网络显著崩溃比例） \\
    \midrule
    Shenyang  & 0.1847 & 0.26 \\
    Zhengzhou & 0.1664 & 0.20 \\
    Chengdu   & 0.1664 & 0.20 \\
    Quanzhou  & 0.1439 & 0.21 \\
    Dongguan  & 0.1458 & 0.21 \\
    Qingdao   & 0.1400 & 0.19 \\
    Harbin    & 0.1360 & 0.06 \\
    Dalian    & 0.1129 & 0.11 \\
    \bottomrule
  \end{tabular}
  \begin{minipage}{\linewidth}
    \vspace{2pt}
    \footnotesize 注：表中城市按 \(R_{bc}\) 降序排列，\(R_{bc}\) 越大表示网络整体鲁棒性越强；
    \(q\) 为随机故障下网络发生显著崩溃时对应的节点移除比例。
  \end{minipage}
\end{table}

\subsection{逐点删除优化（第三题）不同算法效率测量}

我们测量了4.4.1中提到的不同启发优先级算法对于全部城市求到的下降过程。

对于原指标$P(q)=\frac{|LCC(G')|}{N_0}$,如图11及表3：
\begin{figure}[H]
    \centering
    \includegraphics[width=0.95\textwidth]{figures/Q3/Q3_All_Attacks_Pq.pdf}
    \caption{原指标下不同算法效果测量}
\end{figure}

\begin{table}[htbp]
  \centering
  \caption{各城市不同攻击策略下的鲁棒性指标 $R_{P_q}$（LCC）}
  \label{tab:Rpq_all}
  \begin{tabular}{lcccccc}
    \toprule
    \textbf{城市} & \textbf{Degree} & \textbf{Betweenness} & \textbf{Closeness}
      & \textbf{CI\_Radius2} & \textbf{CoreHD} & \textbf{CABS} \\
    \midrule
    Chengdu   & 0.1475 & 0.1240 & 0.3545 & 0.1755 & 0.1131 & \textbf{0.0411} \\
    Dalian    & 0.0989 & 0.0652 & 0.3043 & 0.1245 & 0.0903 & \textbf{0.0336} \\
    Dongguan  & 0.1259 & 0.0957 & 0.2710 & 0.1256 & 0.0942 & \textbf{0.0393} \\
    Harbin    & 0.1393 & 0.1116 & 0.2750 & 0.1678 & 0.1055 & \textbf{0.0385} \\
    Qingdao   & 0.1465 & 0.0735 & 0.3139 & 0.1426 & 0.1079 & \textbf{0.0375} \\
    Quanzhou  & 0.1447 & 0.1141 & 0.2463 & 0.1695 & 0.0965 & \textbf{0.0468} \\
    Shenyang  & 0.1500 & 0.1210 & 0.3037 & 0.1785 & 0.1327 & \textbf{0.0524} \\
    Zhengzhou & 0.1542 & 0.1173 & 0.2248 & 0.2117 & 0.1329 & \textbf{0.0479} \\
    \bottomrule
  \end{tabular}
  \begin{minipage}{\linewidth}
    \vspace{2pt}
    \footnotesize 注：CABS 列对应参数配置 W3\_K30\_B5；\(R\) 值越小表示攻击效果越强。
  \end{minipage}
\end{table}

而对于强化指标$Q_{bc}$,如图12及表4：
\begin{figure}[H]
    \centering
    \includegraphics[width=0.95\textwidth]{figures/Q3/Q3_All_Attacks_Qbc.pdf}
    \caption{原指标下不同算法效果测量}
\end{figure}

\begin{table}[htbp]
  \centering
  \caption{各城市不同攻击策略下的鲁棒性指标 $R_{Q_{bc}}$}
  \label{tab:Rqbc_all}
  \begin{tabular}{lcccccc}
    \toprule
    \textbf{城市} & \textbf{Degree} & \textbf{Betweenness} & \textbf{Closeness}
      & \textbf{CI\_Radius2} & \textbf{CoreHD} & \textbf{CABS} \\
    \midrule
    Chengdu   & 0.1279 & 0.1012 & 0.3174 & 0.1433 & 0.0897 & \textbf{0.0343} \\
    Dalian    & 0.0829 & 0.0513 & 0.2529 & 0.0961 & 0.0653 & \textbf{0.0243} \\
    Dongguan  & 0.1092 & 0.0773 & 0.2377 & 0.0974 & 0.0714 & \textbf{0.0301} \\
    Harbin    & 0.1111 & 0.0928 & 0.2132 & 0.1291 & 0.0799 & \textbf{0.0293} \\
    Qingdao   & 0.1258 & 0.0575 & 0.2738 & 0.1133 & 0.0799 & \textbf{0.0285} \\
    Quanzhou  & 0.1194 & 0.0962 & 0.2101 & 0.1312 & 0.0725 & \textbf{0.0332} \\
    Shenyang  & 0.1301 & 0.1008 & 0.2737 & 0.1486 & 0.1057 & \textbf{0.0406} \\
    Zhengzhou & 0.1313 & 0.0969 & 0.1921 & 0.1703 & 0.1005 & \textbf{0.0396} \\
    \bottomrule
  \end{tabular}
  \begin{minipage}{\linewidth}
    \vspace{2pt}
    \footnotesize 注：CABS 列对应参数配置 W3\_K30\_B5；\(R\) 值越小表示攻击效果越强。
  \end{minipage}
\end{table}

\subsection{逐点删除健壮性测量}
如表5。在两种指标下，均呈现了同样的健壮性排名。
\begin{table}[htbp]
  \centering
  \caption{各城市网络鲁棒性排名（CABS 攻击策略，\(R\) 值越大表示网络越难瓦解）}
  \label{tab:robustness_ranking}
  \begin{tabular}{clcclc}
    \toprule
    \multicolumn{3}{c}{\(R_{P_q}\)（LCC）} & \multicolumn{3}{c}{\(R_{Q_{bc}}\)} \\
    \cmidrule(lr){1-3} \cmidrule(lr){4-6}
    \textbf{排名} & \textbf{城市} & \(R_{P_q}\) &
    \textbf{排名} & \textbf{城市} & \(R_{Q_{bc}}\) \\
    \midrule
    1 & Shenyang  & 0.0524 & 1 & Shenyang  & 0.0406 \\
    2 & Zhengzhou & 0.0479 & 2 & Zhengzhou & 0.0396 \\
    3 & Quanzhou  & 0.0468 & 3 & Chengdu   & 0.0343 \\
    4 & Chengdu   & 0.0411 & 4 & Quanzhou  & 0.0332 \\
    5 & Dongguan  & 0.0393 & 5 & Dongguan  & 0.0301 \\
    6 & Harbin    & 0.0385 & 6 & Harbin    & 0.0293 \\
    7 & Qingdao   & 0.0375 & 7 & Qingdao   & 0.0285 \\
    8 & Dalian    & 0.0336 & 8 & Dalian    & 0.0243 \\
    \bottomrule
  \end{tabular}
  \begin{minipage}{\linewidth}
    \vspace{2pt}
    \footnotesize 注：两项指标均采用 CABS（W3\_K30\_B5）攻击序列计算；
    排名第 1 表示该城市网络结构最为鲁棒。
  \end{minipage}
\end{table}

\subsection{空间级联失效下不同算法的收益分析}

在空间级联失效模型（即当某个节点遭到破坏时，其周边半径$r$内的所有节点均会受到波及而失效）下，我们对比了传统的基于节点度的空间破坏策略（S-Degree）、文献前沿的S-CoreHD策略，以及本文提出的区域集束搜索策略（RCABS-BE）对系统健壮性的影响。图\ref{fig:Q4_R_Combined}展示了在不同破坏半径下，8个城市网络的健壮性积分$R$（以$R_{P_q}$为例）的变化趋势。

\begin{figure}[H]
    \centering
    \includegraphics[width=0.95\textwidth]{figures/Q4/Q4_MultiRadius_RIntegral_Combined.pdf}  
    \caption{不同空间攻击策略横跨不同破坏半径时的$R$积分对比}
    \label{fig:Q4_R_Combined}
\end{figure}

从图\ref{fig:Q4_R_Combined}中可以明显看出，随着影响半径$r$的增加，三种策略造成的健壮性积分$R$均快速下降，这表明大范围的级联瘫痪效应极大地恶化了网络整体的抗毁毁损能力。横向对比来看，本文基于局部致密性与桥接属性优化所提出的 RCABS-BE 算法（下方的橙色曲线）始终保持了最为显著的瓦解压制力。这说明通过动态计算区域网络密度以及基于双连通性质($Q_{bc}$)构建启发式的奖励代价，RCABS-BE能够以极为限缩的攻击步数，优先摧毁城市路网中的“枢纽连接带”，实现“以点破面”的降维打击收益，表现显著优于常规解法。

\subsection{基于RCABS-BE算法的空间健壮性分析（区域瘫痪探讨）}

为了更细致地观察波及半径对各个城市的物理连通性的动态影响，我们在图\ref{fig:Q4_PQ_RCABS}中描绘了由RCABS-BE策略执行时，各城市的最大连通分量占比$P(q)$随被移除节点比例$q$的级联衰减曲线。

\begin{figure}[H]
    \centering
    \includegraphics[width=0.95\textwidth]{figures/Q4/Q4_MultiRadius_Pq_RCABS-BE.pdf}  
    \caption{RCABS-BE策略下不同破坏半径时的$P(q)$级联衰减过程}
    \label{fig:Q4_PQ_RCABS}
\end{figure}

如图\ref{fig:Q4_PQ_RCABS}所示，波及半径$r$起到了决定性的作用。在小尺度的破坏半径（如100m, 250m等浅色与冷色线条）下，网络的结构崩溃呈现出较为平滑的下降序列，城市路网凭借其自身的冗余分支结构勉强维持了部分的物理畅通。然而，一旦破坏半径的烈度超过了某些城市的网络容错阈值（如超过1000m、2000m的暖色与红线级），$P(q)$曲线往往在攻击序列初期仅去除了不到$5\%$的节点时，便发生了不可逆的断崖式跌落。这标志着此时少数几个高度核心的片区（例如中心城区商业核心与跨江桥梁）的连锁瘫痪，瞬间引发了城市交通网络的“全域雪崩”拓扑割裂反应。

不同城市的表现进一步展现了其网络规划的先天差异。例如，郑州或大连的部分曲线在深层阶段仍存留长长的断尾，证明其次级路网分摊了冲击；而在连通形式较为脆弱的中心极化城市中，核心枢纽被锁定后的切除，会导致区域通行能力瞬间化为无数极小的孤岛。

表\ref{tab:Q4_spatial_robustness_ranking}汇总了在典型的灾害波及半径（$r=1000$m）下，执行 RCABS-BE 空间攻击策略后各个城市的健壮性排名。与随机点故障（第三问）的排名结果相比，由于在空间模型下密集连接的枢纽被“一窝端”，之前能够依靠冗余支路抵抗单点破坏的城市，在面临区域性毁灭打击时，其抗毁排名发生了一定的洗牌现象。

\begin{table}[htbp]
  \centering
  \caption{各城市网络空间鲁棒性排名（RCABS-BE 攻击策略，$r=1000$m，$R$ 值越大表示网络越难瓦解）}
  \label{tab:Q4_spatial_robustness_ranking}
  \begin{tabular}{clcclc}
    \toprule
    \multicolumn{3}{c}{\(R_{P_q}\)（LCC积分）} & \multicolumn{3}{c}{\(R_{Q_{bc}}\)（双连通积分）} \\
    \cmidrule(lr){1-3} \cmidrule(lr){4-6}
    \textbf{排名} & \textbf{城市} & \(R_{P_q}\) &
    \textbf{排名} & \textbf{城市} & \(R_{Q_{bc}}\) \\
    \midrule
    1 & Shenyang  & 0.3096 & 1 & Shenyang  & 0.2750 \\
    2 & Zhengzhou & 0.2858 & 2 & Zhengzhou & 0.2417 \\
    3 & Harbin    & 0.2635 & 3 & Chengdu   & 0.2194 \\
    4 & Chengdu   & 0.2517 & 4 & Harbin    & 0.2156 \\
    5 & Qingdao   & 0.2114 & 5 & Qingdao   & 0.1764 \\
    6 & Dongguan  & 0.2096 & 6 & Dongguan  & 0.1699 \\
    7 & Quanzhou  & 0.1930 & 7 & Quanzhou  & 0.1645 \\
    8 & Dalian    & 0.1706 & 8 & Dalian    & 0.1379 \\
    \bottomrule
  \end{tabular}
  \begin{minipage}{\linewidth}
    \vspace{2pt}
    \footnotesize 注：两项指标均采用 RCABS-BE 空间级联攻击序列（波及半径 $r=1000$m）计算；
    排名第 1 表示在该半径的区域摧毁模式下，该城市网络结构存活能力最强。
  \end{minipage}
\end{table}
